\uuid{WVT9}
\titre{Équations complexes}
\auteur{bodin ; hanani ; mzari}
\organisation{exo7}
\licence{CC BY-NC-SA 4.0}
\datecreate{2020-01-20}
\niveau{L1}
\module{Algèbre}
\chapitre{Nombres complexes}
\sousChapitre{Racine carrée, équation du second degré}
\difficulte{1}
\theme{Équations, équations complexes}
\exo7qcmid{lille-1.nombres-complexes.017}

\contenu{
\begin{question}
\qclassification{104.02,104.03,104.04}
\qtags{motcle=équations complexes}
On considère l'équation : $(E) : \, z^2+z+1=0$, $z\in \Cc$.   Quelles sont les assertions vraies ?
\begin{answers}
    \bad{Les solutions de $(E)$ sont : $z_1= \frac{-1+\sqrt5}{2}$ et $z_2= -\frac{1+\sqrt5}{2}$.}

    \good{Les solutions de $(E)$ sont : $z_1= \frac{-1+i\sqrt3}{2}$ et $z_2= -\frac{1+i\sqrt3}{2}$.}

    \good{Les solutions de $(E)$ sont : $z_1= e^{\frac{2i\pi}{3}}$ et $z_2=e^{\frac{-2i\pi}{3}}$.}

    \good{Si $z$ est une solution de $(E)$, alors $|z|=1$.}
\end{answers}
\begin{explanations}
Les solutions complexes d'une équation du second degré $az^2+bz+c=0$ sont $z_1=\frac{-b+\delta}{2a}$ et
$z_1=\frac{-b-\delta}{2a}$, où $\delta$ est une racine carrée de $\Delta=b^2-4ac$.
\end{explanations}
\end{question}
}
